\section{Precision matrix and partial correlation}
\label{subsec: pcorrandprec}
In this part, we introduce the definition of precision matrix and partial correlation. 
Some useful properties of them are discussed. 
Below we suppose random variables $X=(X_1,X_2,\cdots,X_n)$ are jointly Gaussian with 0 mean and covariance $\Sigma$,  
and denote $X-\{X_i\}$ by $V_i$. 

\begin{definition}(\textbf{Precision matrix})

    The precision matrix $\Theta$ of multivariate Gaussian random variables $(X_1,\cdots,X_n)$ is the inversion of the population covariance $\Sigma$.
\end{definition}

The precision matrix of multivariate Gaussian random variables can be used to test the conditional independence. 
In fact, given values of all other identities, 
any two variables $X_i$ and $X_j$ are conditionally independent if $\Theta_{ij}=0$. 
This is shown in the following theorem.
\begin{theorem}
    \label{thm: ch1_conind}
    $X_i$ and $X_j$ are conditionally independent given $V_{i,j}$ if and only if $\Theta_{ij}=\Theta_{ji}=0$.
\end{theorem}
\begin{proof}
    Without loss of generality we assume $i=1$, $j=2$. For short, we simplify $V_{i,j}$ by $V$.
    The joint Gaussian distribution of $(X_1, X_2)$ condition on $V$ is given by (ignoring the constant) 
    \begin{align*}
        p(X_1,X_2|V)=&|\Theta|^{\frac{1}{2}}\exp\{-\frac{1}{2}[X_1 X_2]\Theta\begin{bmatrix}
            X_1 \\
            X_2 
        \end{bmatrix}\}\\
        =& \text{det}\bigg(\begin{bmatrix}
            \Theta_{11} & \Theta_{12} \\
            \Theta_{21} & \Theta_{22} 
        \end{bmatrix}\bigg)^{\frac{1}{2}}\exp\{-\frac{1}{2}[X_1 X_2]\begin{bmatrix}
            \Theta_{11} & \Theta_{12} \\
            \Theta_{21} & \Theta_{22}
        \end{bmatrix}\begin{bmatrix}
            X_1 \\
            X_2
        \end{bmatrix}\} \\
        =& |\Theta_{11}|^{\frac{1}{2}}\exp\{-\frac{1}{2}X_1\Theta_{11}X_1\}\cdot|\Theta_{22}|^{\frac{1}{2}}\exp\{-\frac{1}{2}X_2\Theta_{22}X_2\} \\
        =& p(X_1|V)\cdotp(X_2|V),
    \end{align*}
    if and only if $\Theta_{12}=\Theta_{21}=0$. Since any $X_i$ and $X_j$ can be moved to the place of $X_1$ and $X_2$ by permutation, this completes the proof. 
\end{proof}

The precision matrix determines the conditional dependence between two Gaussian random variables. 
And the partial correlation measures the relation strength of the two random variables when holding other random variables constant. 
The relation strength is featured by the residuals of $X_i$ and $X_j$ regressing on the rest random variables respectively. 

\begin{definition}(\textbf{Partial correlation})
    
    The partial correlation $\rho_{i,j|V_{i,j}}$ between $X_i$ and $X_j$ given the rest of variables
    is the correlation between the residual resulting from regressing $X_i$ on $V_{i,j}$ and $X_j$ on $V_{i,j}$. 
    More formally, denote residuals by $R_i=X_i-\beta_i^TV_{i,j}$ and $R_j=X_j-\beta^T_jV_{i,j}$. 
    Then, 
    \begin{equation}
        \label{eq: ch1_pcorrdef}
        \rho_{i,j|V_{i,j}} = \frac{\text{cov}(R_i,R_j)}{\sqrt{\text{Var}(R_i)}\sqrt{\text{Var}(R_j)}}.
    \end{equation} 
\end{definition}

Since the residual of $X_i$ regressing on $V_{i,j}$ is the orthogonal projection of the linear space spanned by $X$ on the orthogonal complement of the linear space spanned by $V$, 
it excludes all information that can be explained by the linear combination of $V$, 
whereby it contains the direct linear influence of $X_j$ on $X_i$. 
In addition, the partial correlation has close relationship with the precision matrix. 
\begin{proposition}
    \label{prop: ch1_pcorr}
    Let $\Theta$ be the precision matrix of joint distribution of $X$.
    The partial correlation is equivalent to
    \begin{equation*}
        \rho_{i,j|V_{i,j}}=-\frac{\Theta_{ij}}{\sqrt{\Theta_{ii}\Theta_{jj}}}.
    \end{equation*}
\end{proposition}
\begin{proof}
    Without loss of generality we assume $i=1$, $j=2$. For short, we simplify $V_{i,j}$ by $V$. 
    Since $\beta_1$ and $\beta_2$ are least square solutions for problems
    \begin{align*}
        &\arg\min_{\beta\in\mathbb{R}^{n-2}} \mathbb{E}\|X_1-\beta^TV\|_2^2, \\
        &\arg\min_{\beta\in\mathbb{R}^{n-2}} \mathbb{E}\|X_2-\beta^TV\|_2^2,
    \end{align*}
    and have closed-form $\beta_1=\Sigma_{VV}^{-1}\Sigma_{VX_1}$ and $\beta_2=\Sigma_{VV}^{-1}\Sigma_{VX_2}$ for 0 centered random variables, 
    Equation (\ref{eq: ch1_pcorrdef}) can be written by 
    \begin{align}
        \label{eq: ch1_pcorrprf1}
        \rho_{1,2|V} &= \frac{\text{cov}(R_1,R_2)}{\sqrt{\text{Var}(R_1)}\sqrt{\text{Var}(R_2)}} \notag \\
                     &= \frac{\text{cov}(X_1-\Sigma_{VX_1}^T\Sigma_{VV}^{-1}V, X_2-\Sigma_{VX_2}^T\Sigma_{VV}^{-1}V)}{\sqrt{\text{Var}(X_1-\Sigma_{VX_1}^T\Sigma_{VV}^{-1}V)}\sqrt{\text{Var}(X_2-\Sigma_{VX_2}^T\Sigma_{VV}^{-1}V)}}.
    \end{align}
    For nominator, we have 
    \begin{align*}
        \text{cov}(X_1-\Sigma_{VX_1}^T\Sigma_{VV}^{-1}V, X_2-\Sigma_{VX_2}^T\Sigma_{VV}^{-1}V) &= \mathbb{E}\left[(X_1-\Sigma_{VX_1}^T\Sigma_{VV}^{-1}V)(X_2-\Sigma_{VX_2}^T\Sigma_{VV}^{-1}V)\right] & (\mathbb{E}(R_1)=\mathbb{E}(R_2)=0) \\
                                                                                               &= \Sigma_{X_1X_2} - \Sigma_{VX_2}^{T}\Sigma_{VV}^{-1}\Sigma_{VX_1}.
    \end{align*}
    For denominator, we have 
    \begin{align*}
        \text{Var}(X_1-\Sigma_{VX_1}^T\Sigma_{VV}^{-1}V) &= \mathbb{E}\left[(X_1-\Sigma_{VX_1}^T\Sigma_{VV}^{-1}V)(X_1-\Sigma_{VX_1}^T\Sigma_{VV}^{-1}V)\right] \\
                                                         &= \Sigma_{X_1X_1} -\Sigma_{VX_1}^T\Sigma_{VV}^{-1}\Sigma_{VX_1}.
    \end{align*}
    Hence, Equation (\ref{eq: ch1_pcorrprf1}) equals to 
    \begin{equation}
        \label{eq: ch1_pcorrprf2}
        \frac{\Sigma_{X_1X_2} - \Sigma_{VX_2}^{T}\Sigma_{VV}^{-1}\Sigma_{VX_1}}{\sqrt{\Sigma_{X_1X_1} -\Sigma_{VX_1}^T\Sigma_{VV}^{-1}\Sigma_{VX_1}}\sqrt{\Sigma_{X_2X_2} - \Sigma_{VX_2}^{T}\Sigma_{VV}^{-1}\Sigma_{VX_2}}}.
    \end{equation}
    Let $\Sigma=\begin{bmatrix}
        \Sigma_{X_1X_1} & \Sigma_{X_1X_2} & \Sigma_{X_1V} \\
        \Sigma_{X_2X_1} & \Sigma_{X_2X_2} & \Sigma_{X_2V} \\
        \Sigma_{VX_1} & \Sigma_{VX_2} & \Sigma_{VV}
    \end{bmatrix}=\begin{bmatrix}
        \Sigma_{11} & \Sigma_{12} \\
        \Sigma_{21} & \Sigma_{VV}   
    \end{bmatrix}$. By Schur's formula for block-matrix inversion, 
    $\Theta=\Sigma^{-1}=\begin{bmatrix}
        (\Sigma_{11}-\Sigma_{12}\Sigma_{VV}^{-1}\Sigma_{21})^{-1} & -(\Sigma_{11}-\Sigma_{12}\Sigma_{VV}^{-1}\Sigma_{21})^{-1}\Sigma_{12}\Sigma_{VV}^{-1} \\
        -\Sigma_{VV}^{-1}\Sigma_{21}(\Sigma_{11}-\Sigma_{12}\Sigma_{VV}^{-1}\Sigma_{21})^{-1} & \Sigma_{VV}^{-1}+\Sigma_{VV}^{-1}\Sigma_{21}(\Sigma_{11}-\Sigma_{12}\Sigma_{VV}^{-1}\Sigma_{21})^{-1}\Sigma_{12}\Sigma_{VV}^{-1}
    \end{bmatrix}$.

    With the left-upper corner of $\Theta$ which is
    \begin{align*}
        (\Sigma_{11}-\Sigma_{12}\Sigma_{VV}^{-1}\Sigma_{21})^{-1} &= \bigg(\begin{bmatrix}
            \Sigma_{X_1X_1} & \Sigma_{X_1X_2} \\
            \Sigma_{X_2X_1} & \Sigma_{X_2X_2}
        \end{bmatrix}-\begin{bmatrix}
            \Sigma_{X_1V}\Sigma_{VV}^{-1}\Sigma_{VX_1} & \Sigma_{X_1V}\Sigma_{VV}^{-1}\Sigma_{VX_2} \\
            \Sigma_{X_2V}\Sigma_{VV}^{-1}\Sigma_{VX_1} & \Sigma_{X_2V}\Sigma_{VV}^{-1}\Sigma_{VX_2}
        \end{bmatrix}\bigg)^{-1} \\
                &= \frac{1}{\text{det}}\begin{bmatrix}
                    \Sigma_{X_2X_2}-\Sigma_{X_2V}\Sigma_{VV}^{-1}\Sigma_{VX_2} & -\Sigma_{X_2X_1}+\Sigma_{X_2V}\Sigma_{VV}^{-1}\Sigma_{VX_1} \\
                    -\Sigma_{X_1X_2}+\Sigma_{X_1V}\Sigma_{VV}^{-1}\Sigma_{VX_2} & \Sigma_{X_1X_1}-\Sigma_{X_1V}\Sigma_{VV}^{-1}\Sigma_{VX_1}
                \end{bmatrix} \\
                &= \begin{bmatrix}
                    \Theta_{11} & \Theta_{12} \\
                    \Theta_{21} & \Theta_{22}
                \end{bmatrix},
    \end{align*}
    we observe that Equation (\ref{eq: ch1_pcorrprf2}) is just $-\frac{\Theta_{12}}{\sqrt{\Theta_{11}\Theta_{22}}}$.
    Since any $X_i$ and $X_j$ can be moved to the place of $X_1$ and $X_2$ by permutation, this completes the proof. 
\end{proof}

By the Proposition \ref{prop: ch1_pcorr}, a direct corollary of Theorem \ref{thm: ch1_conind} is obtained. 
\begin{corollary}
    $X_i$ and $X_j$ are conditionally independent given $V_{i,j}$ if and only if $\rho_{i,j|V_{i,j}}=\rho_{j,i|V_{i,j}}=0$.
\end{corollary}